%-------------------------
% Resume in Latex
% Author : Aditya Reddy
% Based off of: https://github.com/sb2nov/resume
% License : MIT
%------------------------
\documentclass[letterpaper,11pt]{article}
\usepackage{latexsym}
\usepackage[empty]{fullpage}
\usepackage{titlesec}
\usepackage{marvosym}
\usepackage[usenames,dvipsnames]{color}
\usepackage{enumitem}
\usepackage[hidelinks]{hyperref}
\usepackage{fancyhdr}
\usepackage[english]{babel}
\usepackage{tabularx}
\input{glyphtounicode}
\pagestyle{fancy}
\fancyhf{}
\fancyfoot{}
\renewcommand{\headrulewidth}{0pt}
\renewcommand{\footrulewidth}{0pt}
\addtolength{\oddsidemargin}{-0.5in}
\addtolength{\evensidemargin}{-0.5in}
\addtolength{\textwidth}{1.0in}
\addtolength{\topmargin}{-0.65in}
\addtolength{\textheight}{1.3in}
\urlstyle{same}
\raggedbottom
\raggedright
\setlength{\tabcolsep}{0in}
\titleformat{\section}{
  \vspace{-5pt}\scshape\raggedright\large
}{}{0em}{}[\color{black}\titlerule\vspace{-5pt}]
\pdfgentounicode=1

\newcommand{\resumeItem}[1]{
  \item\small{{#1\vspace{-1pt}}}
}
\newcommand{\resumeSubheading}[4]{
  \vspace{-1pt}\item
    \begin{tabular*}{0.97\textwidth}[t]{l@{\extracolsep{\fill}}r}
      \textbf{#1} & \textbf{#2} \\
      \textit{\small#3} & \textit{\small #4} \\
    \end{tabular*}\vspace{-7pt}
}

% Two-line project heading: {Name}{Link}{Tools/Stack}
\newcommand{\resumeProjectHeading}[3]{
    \item
    \begin{tabular*}{0.97\textwidth}{l@{\extracolsep{\fill}}r}
      \textbf{\small#1} & \textbf{\small #2} \\
      \textit{\small#3} & \\
    \end{tabular*}\vspace{-7pt}
}

\newcommand{\resumeSubItem}[1]{\resumeItem{#1}\vspace{-4pt}}
\renewcommand\labelitemii{$\vcenter{\hbox{\tiny$\bullet$}}$}
\newcommand{\resumeSubHeadingListStart}{\begin{itemize}[leftmargin=0.0in, label={}]}
\newcommand{\resumeSubHeadingListEnd}{\end{itemize}}
\newcommand{\resumeItemListStart}{\begin{itemize}[leftmargin=0.15in, itemsep=0pt, parsep=0pt, topsep=1pt]}
\newcommand{\resumeItemListEnd}{\end{itemize}\vspace{-2pt}}

\begin{document}

%----------HEADING----------
\begin{center}
    \textbf{\Huge Aditya Reddy} \\ \vspace{3pt}
    \href{mailto:adityareddy400@gmail.com}{\underline{adityareddy400@gmail.com}} $|$
    \href{https://github.com/Aditya-1020}{\underline{github.com/Aditya-1020}} $|$
    \href{https://aditya-1020.github.io}{\underline{aditya-1020.github.io}}
\end{center}

%-----------EDUCATION-----------
\section{Education}
\resumeSubHeadingListStart
  \resumeSubheading
    {Manipal University Jaipur}{Aug 2024 -- Expected May 2028}
    {Bachelor of Technology, Electronics and Communication Engineering}{Jaipur, India}
\resumeSubHeadingListEnd

%-----------TECHNICAL SKILLS-----------
\section{Technical Skills}
\begin{itemize}[leftmargin=0.0in, label={}, itemsep=1pt]
  \item\small{\textbf{HDL / RTL Design:} SystemVerilog, Verilog}
  \item\small{\textbf{Physical Design:} RTL-to-GDS flow, Floorplanning, Placement \& Route (PnR), Clock Tree Synthesis (CTS), Static Timing Analysis (STA), DRC/LVS signoff}
  \item\small{\textbf{EDA \& Tools:} OpenLane~2, OpenROAD, Yosys, Vivado, Magic, KLayout, OpenSTA, SymbiYosys, Verilator, GTKWave, LTspice; SKY130 PDK}
  \item\small{\textbf{Programming:} C, C++, Python}
  \item\small{\textbf{Embedded:} ESP32, Raspberry Pi, Arduino}
  \item\small{\textbf{Systems:} Linux, Git, Make, GCC, RISC-V toolchain, POSIX}
\end{itemize}

%-----------PROJECTS-----------
\section{Projects}
\resumeSubHeadingListStart

  \resumeProjectHeading
    {Systolic Array Accelerator -- 4$\times$4 Output-Stationary MAC Array}
    {\href{https://github.com/Aditya-1020/AdiSystolic}{GitHub}}
    {SystemVerilog, OpenLane~2, SKY130B PDK}
  \resumeItemListStart
    \resumeItem{Designed a 4$\times$4 systolic array with 16 pipelined Multiply-Accumulate (MAC) processing elements (PEs) and a ping-pong dual-bank register file to overlap BRAM load latency with compute, eliminating inter-matrix stall cycles; validated all 48 FC2 inference tiles cycle-accurate against a golden INT8 MNIST model.}
    \resumeItem{Ran full RTL-to-GDS on SKY130B via OpenLane~2, diagnosing and resolving hold violations through targeted post-placement buffer insertion; achieved Static Timing Analysis (STA) closure at 76~MHz worst-corner (SS/125\,\textdegree{}C) and 137~MHz nominal (TT/27\,\textdegree{}C) with zero setup, hold, Design Rule Check (DRC), and Layout vs.\ Schematic (LVS) violations across all 9 PVT corners; 18\,K cells, 0.21\,mm$^{2}$ core area.}
  \resumeItemListEnd

  \vspace{0pt}

  \resumeProjectHeading
    {RISC-V Processor Core -- RV32IM 5-Stage Pipeline}
    {\href{https://github.com/Aditya-1020/AdiRiscV}{GitHub}}
    {SystemVerilog, Vivado, Verilator}
  \resumeItemListStart
    \resumeItem{Implemented a full 5-stage RV32IM pipeline with EX/MEM forwarding, load-use stall insertion, and 28 SystemVerilog Assertions (SVA) at 100\% pass rate across 8 test programs; added single-cycle 64-bit hardware multiply and a 32-cycle non-restoring radix-2 divider that stalls the pipeline via a \texttt{ready} signal, covering all M-extension edge cases.}
    \resumeItem{Built a 3-component dynamic branch predictor -- 16-entry Branch Target Buffer (BTB), 256-entry Gshare Pattern History Table (PHT), and 4-deep Return Address Stack (RAS) -- synthesised to Artix-7 in Vivado, closing timing at 100~MHz with performance counters quantifying IPC, branch mispredict rate, and stall cycles.}
  \resumeItemListEnd

  \vspace{0pt}

\resumeProjectHeading
    {Dual-Clock Asynchronous FIFO -- Formally Verified CDC Design}
    {\href{https://github.com/Aditya-1020/async}{GitHub}}
    {Verilog, OpenRAM, SymbiYosys, cocotb}
  \resumeItemListStart
    \resumeItem{Designed a parametric dual-clock FIFO in Verilog with Gray-coded read/write pointers and 2-flop Clock Domain Crossing (CDC) synchronizers; instantiated an OpenRAM pseudo-dual-port SRAM macro (16$\times$64, SKY130) for the storage array, swapped for a combinational stub during formal to bypass non-synthesizable timing constructs.}
    \resumeItem{Formally verified no-overflow, no-underflow, and CDC safety using SymbiYosys with smtbmc/Boolector --- BMC at depth 128 with \texttt{multiclock on} for the top-level, and k-induction (\texttt{mode prove}) for \texttt{rst\_sync}, \texttt{fifo\_mem}, and \texttt{sync\_fifo} sub-modules; validated with 13 cocotb testcases across \texttt{async\_fifo} and \texttt{fifo\_mem}, covering reset state, full/empty conditions, pointer wraparound, concurrent read/write, and asymmetric clock ratios up to 4:1.}
  \resumeItemListEnd

  \vspace{0pt}

  \resumeProjectHeading
    {USB-C Supply Voltage Monitor}
    {}
    {Analog Circuit Design, LTspice}
  \resumeItemListStart
    \resumeItem{Designed a multi-threshold LM324 op-amp comparator network to classify USB-C rail voltage into three bands ($<$4\,V, 4--5\,V, $>$5\,V) with LED indicators; verified hysteresis-free switching across the full DC sweep in LTspice before confirming on a breadboarded prototype.}
  \resumeItemListEnd

\resumeSubHeadingListEnd

\vspace{0pt}

%-----------ACTIVITIES-----------
\section{Activities}
\begin{itemize}[leftmargin=0in, label={}, itemsep=0pt]
  \item\small{\textbf{Volunteer} --- Bison Asha School for Special Needs: supported educational activities and organised inclusive events \hfill 2022--2024}
\end{itemize}

\end{document}